\chapter{Le liste collegate}
\label{chap:liste}

Gli array hanno una dimensione fissa, e i loro elementi stanno uno accanto
all'altro: per far posto a un elemento nel mezzo bisogna spostare tutti
quelli che seguono. Le \textbf{liste collegate} superano questi limiti con
un'idea semplice: ogni elemento sa dove si trova il successivo. In questo
capitolo mettiamo insieme ciò che abbiamo visto finora (strutture, puntatori,
\code{malloc} e \code{free}) per costruire una lista e scrivere le operazioni
elencate nelle slide; poi introduciamo la \textbf{ricorsione} e i
\textbf{file}, e con questi strumenti leggiamo per intero il progetto delle
lezioni incontrato nel capitolo~\ref{chap:make}.

\section{Perché le liste}

\subsection{I limiti degli array}

Un array ha due limiti. Il primo: la sua dimensione si sceglie quando lo si
crea; se poi i dati crescono, bisogna allocare un array più grande e
ricopiarli. Il secondo: gli elementi sono \textbf{consecutivi}, quindi per
inserire un elemento in mezzo bisogna spostare di un posto tutti quelli che
seguono, e per toglierne uno bisogna spostarli indietro
(figura~\ref{fig:array-lista}, in alto). In un array di un milione di
elementi, inserirne uno all'inizio richiede un milione di spostamenti.

\subsection{L'idea: ogni elemento sa dov'è il successivo}

Una \textbf{lista collegata} (\textit{linked list}) è una struttura dati
\textbf{dinamica} in cui ogni elemento, detto \textbf{nodo}, è collegato al
successivo. Ogni nodo contiene

\begin{itemize}
  \item i \textbf{dati};
  \item un \textbf{puntatore al nodo successivo}.
\end{itemize}

I nodi possono stare ovunque in memoria: l'ordine della lista non dipende da
dove si trovano, ma dai puntatori. Per inserire un elemento basta creare un
nodo e cambiare due puntatori, senza spostare niente
(figura~\ref{fig:array-lista}, in basso).

\begin{figure}[H]
  \centering
  \begin{tikzpicture}[font=\Lato\scriptsize]
    % l'array
    \node[mnota, anchor=west, font=\Lato\scriptsize\bfseries] at (-1.0,1.25)
      {array: per inserire 30 in posizione 2 si spostano 50 e 40};
    \celle{a}{(0,0)}{10,20,40,50,{},{}}
    \draw[mcerca] (a-3.north) to[bend left=70] node[mpasso, above=1pt] {1} (a-4.north);
    \draw[mcerca] (a-2.north) to[bend left=70] node[mpasso, above=1pt] {2} (a-3.north);
    \node[mvar, minimum width=0.6cm, fill=accent!10!white] (x) at (-0.6,0.75) {30};
    \draw[mcerca] (x.east) to[bend left=20] node[mpasso, above=1pt] {3} (a-2.north west);
    \node[mnota, anchor=west] at ([xshift=0.25cm]a-5.east) {prima};
    \celle{b}{(0,-1.15)}{10,20,30,40,50,{}}
    \node[mnota, anchor=west] at ([xshift=0.25cm]b-5.east) {dopo};
    % la lista
    \node[mnota, anchor=west, font=\Lato\scriptsize\bfseries] at (-1.0,-2.3)
      {lista: si crea un nodo e si cambiano due puntatori};
    \variabile[minimum width=0.7cm]{h}{(0,-3.1)}{head}{}
    \nodo{n1}{(0.8,-3.1)}{10}
    \nodo{n2}{(2.7,-3.1)}{20}
    \nodo{n4}{(6.0,-3.1)}{40}
    \nodo{n5}{(7.9,-3.1)}{50}
    \nodonull{n5}
    \nodo[mlnuovo]{n3}{(4.3,-4.25)}{30}
    \riferimento{h}{-- (n1.west)}
    \collega{n1}{n2}
    \collega{n4}{n5}
    \draw[mrifvecchio] (n2-next.center) -- (n4.west);
    \draw[mrifnuovo] (n3-next.center) -- node[mpasso, pos=0.45, right=3pt] {1} (n4.west);
    \draw[mrifnuovo] (n2-next.center) -- node[mpasso, pos=0.55, left=3pt] {2} (n3.west);
    \node[mpunto] at (n2-next.center) {};
    \node[mpunto] at (n3-next.center) {};
  \end{tikzpicture}
  \caption{Inserire 30 fra 20 e 40. In un array (in alto) gli elementi
           successivi si spostano di un posto, a partire dall'ultimo. In una
           lista (in basso) il nuovo nodo viene collegato cambiando due
           puntatori, nell'ordine indicato; la freccia tratteggiata è il
           collegamento che non c'è più.}
  \label{fig:array-lista}
\end{figure}

\subsection{Vantaggi e svantaggi}

\begin{center}
\renewcommand{\arraystretch}{1.15}
\begin{tabular}{@{}lll@{}}
  \toprule
                                 & \textbf{array}              & \textbf{lista} \\
  \midrule
  dimensione                     & fissata quando lo si crea   & cresce e cala un nodo alla volta \\
  inserire o togliere in testa   & si spostano tutti gli elementi & si cambiano due puntatori \\
  arrivare all'elemento $i$      & immediato: \code{a[i]}      & bisogna attraversare $i$ nodi \\
  memoria                        & solo i dati                 & i dati e un puntatore per nodo \\
  \bottomrule
\end{tabular}
\end{center}

Il vantaggio delle liste è la \textbf{flessibilità}: la dimensione cambia a
ogni inserimento, e inserimenti e cancellazioni sono efficienti perché non
spostano dati. Lo svantaggio è l'\textbf{accesso sequenziale}: non si può
saltare direttamente all'elemento $i$, perché il suo indirizzo si scopre
solo arrivandoci, seguendo i puntatori a partire dal primo nodo. In un
array, invece, l'indirizzo di \code{a[i]} si calcola con una moltiplicazione
e una somma (capitolo~\ref{chap:puntatori}).

\section{Tipi di liste}

Esistono diverse varianti (figura~\ref{fig:tipi-liste}):

\begin{itemize}
  \item nelle liste \textbf{semplicemente collegate} (\textit{singly linked
        list}) ogni nodo punta al successivo. Sono quelle di questo capitolo;
  \item nelle liste \textbf{doppiamente collegate} (\textit{doubly linked
        list}) ogni nodo punta al successivo e al precedente, così la lista
        si può percorrere nei due versi (al prezzo di un puntatore in più per
        nodo, da tenere aggiornato);
  \item nelle liste \textbf{circolari} l'ultimo nodo punta al primo,
        formando un ciclo: non c'è un \code{NULL} che segna la fine, e lo
        scorrimento si ferma quando si torna al nodo da cui si è partiti.
\end{itemize}

\begin{figure}[H]
  \centering
  \begin{tikzpicture}[font=\Lato\scriptsize]
    % semplicemente collegata
    \node[mnota, anchor=west, font=\Lato\scriptsize\bfseries] at (-1.0,0.7)
      {semplicemente collegata: ogni nodo punta al successivo};
    \variabile[minimum width=0.7cm]{h1}{(0,0)}{head}{}
    \nodo{s1}{(0.8,0)}{10}
    \nodo{s2}{(2.7,0)}{20}
    \nodo{s3}{(4.6,0)}{30}
    \nodonull{s3}
    \riferimento{h1}{-- (s1.west)}
    \collega{s1}{s2}
    \collega{s2}{s3}
    % doppiamente collegata
    \node[mnota, anchor=west, font=\Lato\scriptsize\bfseries] at (-1.0,-1.0)
      {doppiamente collegata: ogni nodo punta al successivo e al precedente};
    \variabile[minimum width=0.7cm]{h2}{(0,-1.75)}{head}{}
    \foreach \k/\ix/\val in {1/0.8/10, 2/3.45/20, 3/6.1/30} {
      \node[mlnodo, minimum width=0.45cm, anchor=west] (dp\k) at (\ix,-1.75) {};
      \node[mlnodo, minimum width=0.9cm, anchor=west] (dd\k) at ([xshift=-0.7pt]dp\k.east) {\val};
      \node[mlnodo, minimum width=0.45cm, anchor=west] (dn\k) at ([xshift=-0.7pt]dd\k.east) {};
    }
    \node[mindir, anchor=south] at ([yshift=1pt]dp1.north) {prev};
    \node[mindir, anchor=south] at ([yshift=1pt]dd1.north) {data};
    \node[mindir, anchor=south] at ([yshift=1pt]dn1.north) {next};
    \riferimento{h2}{-- (dp1.west)}
    \foreach \k/\j in {1/2, 2/3} {
      \draw[mrif] ([yshift=0.09cm]dn\k.center) -- ([yshift=0.09cm]dp\j.west);
      \node[mpunto] at ([yshift=0.09cm]dn\k.center) {};
      \draw[mrif] ([yshift=-0.09cm]dp\j.center) -- ([yshift=-0.09cm]dn\k.east);
      \node[mpunto] at ([yshift=-0.09cm]dp\j.center) {};
    }
    \foreach \c in {dp1, dn3}
      \draw[line width=0.7pt, draw=primary!65!black]
        ([shift={(1.5pt,1.5pt)}]\c.south west) -- ([shift={(-1.5pt,-1.5pt)}]\c.north east);
    % circolare
    \node[mnota, anchor=west, font=\Lato\scriptsize\bfseries] at (-1.0,-2.75)
      {circolare: l'ultimo nodo punta al primo};
    \variabile[minimum width=0.7cm]{h3}{(0,-3.45)}{head}{}
    \nodo{c1}{(0.8,-3.45)}{10}
    \nodo{c2}{(2.7,-3.45)}{20}
    \nodo{c3}{(4.6,-3.45)}{30}
    \riferimento{h3}{-- (c1.west)}
    \collega{c1}{c2}
    \collega{c2}{c3}
    \draw[mrif] (c3-next.center) -- ++(0.55,0) |- ([yshift=-0.4cm]c1-dato.south) -- (c1-dato.south);
    \node[mpunto] at (c3-next.center) {};
  \end{tikzpicture}
  \caption{Le tre varianti delle liste. La barra in un campo indica il
           puntatore \code{NULL}.}
  \label{fig:tipi-liste}
\end{figure}

In C le varianti differiscono solo per i campi del nodo e per il modo in cui
si riconosce la fine della lista. Il nodo di una lista doppiamente collegata,
per esempio, ha due puntatori:

\begin{lstlisting}[style=c]
struct DNode {
   int data;
   struct DNode *prev;     /* il nodo precedente */
   struct DNode *next;     /* il nodo successivo */
};
\end{lstlisting}

\section{Una lista in C}

\subsection{Il nodo: una struttura che punta a se stessa}

Il nodo di una lista di interi, come nelle slide, è una struttura con due
campi: il dato e il puntatore al nodo successivo.

\begin{lstlisting}[style=c]
struct Node {
   int data;               /* il dato */
   struct Node *next;      /* il puntatore al nodo successivo */
};
\end{lstlisting}

Il campo \code{next} è un puntatore a \code{struct Node}, cioè al tipo che si
sta definendo: una struttura come questa si dice
\textbf{autoreferenziale}. È lecito perché un puntatore occupa sempre lo
stesso spazio (8 byte), qualunque cosa punti: il compilatore sa quanto posto
riservare a \code{next} anche se \code{struct Node} non è ancora completa.
(Le slide scrivono \code{struct Node* next}: lo spazio prima o dopo
l'asterisco non conta.) Non si può invece mettere dentro un nodo un
\textbf{nodo intero}: conterrebbe un nodo, che conterrebbe un nodo, e così
via all'infinito. Il compilatore lo rifiuta, perché mentre la definizione è
in corso il tipo è ancora incompleto:

\begin{console}
$ gcc -Wall -o nodo_intero nodo_intero.c
nodo_intero.c:3:16: error: field (*@‘@*)next(*@’@*) has incomplete type
    3 |    struct Node next;      /* un nodo intero, non un puntatore */
      |                ^~~~
\end{console}

\subsection{La lista vuota}

Una lista si rappresenta con un puntatore al suo primo nodo, che per
tradizione si chiama \code{head} (testa). La lista vuota non ha nodi, e il
puntatore vale \code{NULL}; le slide la creano così:

\begin{lstlisting}[style=c]
struct Node *head = NULL;      /* la lista vuota */
\end{lstlisting}

L'inizializzazione è essenziale: un puntatore non inizializzato contiene un
indirizzo casuale, e il programma lo scambierebbe per il primo nodo di una
lista.

\subsection{Collegare tre nodi a mano}

Per vedere com'è fatta una lista, costruiamone una di tre nodi a mano:
li allochiamo con \code{malloc}, ne riempiamo i dati e li colleghiamo uno
per uno. Poi scorriamo la lista stampando i dati, e stampiamo anche
l'indirizzo di ogni nodo e il valore del suo campo \code{next}:

\begin{affianco}[0.62][c]
\begin{lstlisting}[style=c, name=nodi.c]
#include <stdio.h>
#include <stdlib.h>

struct Node {
   int data;
   struct Node *next;
};

int main(void)
{
   struct Node *head = NULL;
   struct Node *a, *b, *c, *p;

   a = malloc(sizeof(struct Node));
   b = malloc(sizeof(struct Node));
   c = malloc(sizeof(struct Node));
   if (a == NULL || b == NULL || c == NULL)
      return 1;
   a->data = 10;
   b->data = 20;
   c->data = 30;
   a->next = b;            /* dopo a viene b */
   b->next = c;            /* dopo b viene c */
   c->next = NULL;         /* c è l'ultimo */
   head = a;               /* la lista inizia da a */

   for (p = head; p != NULL; p = p->next)
      printf("%d -> ", p->data);
   printf("NULL\n");
   printf("head = %p\n", (void *) head);
   for (p = head; p != NULL; p = p->next)
      printf("nodo %p: data = %d, next = %p\n",
             (void *) p, p->data, (void *) p->next);
   free(a);
   free(b);
   free(c);
   return 0;
}
\end{lstlisting}
\risultato
Ogni nodo è una \code{struct Node} allocata nello heap; \code{a}, \code{b} e
\code{c} sono puntatori ai tre nodi, e servono solo a costruire la lista e a
liberarla alla fine. Le tre assegnazioni ai campi \code{next} creano i
collegamenti; \code{head = a} dice dove comincia la lista. I due cicli
\code{for} scorrono la lista dall'inizio alla fine, come vedremo fra poco.
\end{affianco}

\begin{console}
$ gcc -Wall -o nodi nodi.c
$ ./nodi
10 -> 20 -> 30 -> NULL
head = 0x5608b5fd32a0
nodo 0x5608b5fd32a0: data = 10, next = 0x5608b5fd32c0
nodo 0x5608b5fd32c0: data = 20, next = 0x5608b5fd32e0
nodo 0x5608b5fd32e0: data = 30, next = (nil)
\end{console}

\subsection{Che cosa mostra l'output}

\begin{figure}[H]
  \centering
  \begin{tikzpicture}[font=\Lato\scriptsize,
      dato/.style={mlnodo, minimum width=0.75cm},
      ind/.style={mlnodo, minimum width=2.25cm, font=\FiraCodeNL\tiny}]
    \node[mvar, minimum width=2.25cm, font=\FiraCodeNL\tiny] (h) at (0,0) {0x5608b5fd32a0};
    \node[mnome, anchor=east] at ([xshift=-2pt]h.west) {head};
    \foreach \k/\ix/\val/\nx/\ad/\nome in {1/1.85/10/0x5608b5fd32c0/0x5608b5fd32a0/a,
        2/5.55/20/0x5608b5fd32e0/0x5608b5fd32c0/b, 3/9.25/30/{(nil)}/0x5608b5fd32e0/c} {
      \node[dato, anchor=west] (d\k) at (\ix,0) {\val};
      \node[ind, anchor=west] (x\k) at ([xshift=-0.7pt]d\k.east) {\nx};
      \node[mnota, anchor=south, align=center, inner sep=1pt]
        at ([yshift=2pt]$(d\k.north west)!0.5!(x\k.north east)$)
        {nodo \code{*\nome}, all'indirizzo\\{\FiraCodeNL\tiny\color{neutral!40!black}\ad}};
      \node[mindir, anchor=north] at ([yshift=-1pt]d\k.south) {data};
      \node[mindir, anchor=north] at ([yshift=-1pt]x\k.south) {next};
    }
    \draw[mrif] (h.east) -- (d1.west);
    \draw[mrif] (x1.east) -- (d2.west);
    \draw[mrif] (x2.east) -- (d3.west);
  \end{tikzpicture}
  \caption{La lista di \code{nodi.c} con gli indirizzi reali dell'esecuzione:
           ogni campo \code{next} contiene l'indirizzo del nodo successivo, e
           le frecce sono solo un modo di disegnarlo.}
  \label{fig:nodi}
\end{figure}

\begin{itemize}
  \item \code{head} contiene l'indirizzo del primo nodo, e il campo
        \code{next} di ogni nodo contiene l'\textbf{indirizzo} del nodo
        successivo (figura~\ref{fig:nodi}). Le frecce con cui disegneremo
        le liste rappresentano proprio questi indirizzi.
  \item Il \code{next} dell'ultimo nodo vale \code{NULL}, che \code{printf}
        con \code{\%p} stampa come \code{(nil)}.
  \item Nella prova i nodi distano 32 byte l'uno dall'altro (\code{0x20} in
        esadecimale): 16 sono del nodo, gli altri sono usati da
        \code{malloc} per gestire i blocchi. Sono vicini perché sono stati
        allocati uno dopo l'altro, ma nulla lo garantisce: in un programma
        che inserisce e cancella, i nodi finiscono sparsi nello heap, e
        l'unico modo per sapere dov'è il successivo è leggere \code{next}.
\end{itemize}

\subsection{Scorrere una lista}

Il ciclo che stampa i dati è lo schema di quasi tutte le operazioni sulle
liste, e conviene capirlo bene:

\begin{lstlisting}[style=c]
for (p = head; p != NULL; p = p->next)
   printf("%d -> ", p->data);
\end{lstlisting}

\begin{itemize}
  \item \code{p = head}: un puntatore di appoggio parte dal primo nodo. Si
        usa \code{p} e non \code{head} perché \code{head} non deve cambiare:
        è l'unico modo di ritrovare la lista.
  \item \code{p != NULL}: finché \code{p} punta a un nodo, il corpo del ciclo
        lo elabora (qui ne stampa il dato, \code{p->data}).
  \item \code{p = p->next}: \code{p} passa al nodo successivo, prendendone
        l'indirizzo dal nodo corrente. Dopo l'ultimo nodo \code{p} diventa
        \code{NULL}, e il ciclo termina (figura~\ref{fig:scorri}).
\end{itemize}

\begin{figure}[H]
  \centering
  \begin{tikzpicture}[font=\Lato\scriptsize]
    \variabile[minimum width=0.7cm]{h}{(0,0)}{head}{}
    \nodo{n1}{(0.8,0)}{10}
    \nodo{n2}{(3.0,0)}{20}
    \nodo{n3}{(5.2,0)}{30}
    \nodonull{n3}
    \riferimento{h}{-- (n1.west)}
    \collega{n1}{n2}
    \collega{n2}{n3}
    \foreach \k/\testo in {1/{primo giro:\\stampa 10}, 2/{secondo giro:\\stampa 20},
        3/{terzo giro:\\stampa 30}} {
      \node[mvar, minimum width=0.6cm] (p\k) at ([yshift=-1.0cm]n\k-dato.south) {};
      \node[mnome, anchor=east] at ([xshift=-2pt]p\k.west) {p};
      \riferimento{p\k}{-- (n\k-dato.south)}
      \node[mnota, anchor=north, align=center] at ([yshift=-2pt]p\k.south) {\testo};
    }
    \node[mvar, minimum width=0.9cm, font=\FiraCodeNL\tiny] (p4) at ([xshift=2.2cm]p3) {NULL};
    \node[mnome, anchor=east] at ([xshift=-2pt]p4.west) {p};
    \node[mnota, anchor=north, align=center] at ([yshift=-2pt]p4.south) {\code{p == NULL}:\\il ciclo finisce};
  \end{tikzpicture}
  \caption{Il puntatore \code{p} durante lo scorrimento: a ogni giro prende
           l'indirizzo contenuto nel campo \code{next} del nodo su cui si
           trova.}
  \label{fig:scorri}
\end{figure}

Lo stesso ciclo, con un contatore al posto della \code{printf}, calcola la
\textbf{lunghezza} della lista. Nel modulo che scriveremo fra poco diventano
due funzioni:

\begin{affianco}[0.5][c]
\begin{lstlisting}[style=c, name=lista.c]
void stampaLista(struct Node *head)
{
   struct Node *p;

   for (p = head; p != NULL; p = p->next)
      printf("%d -> ", p->data);
   printf("NULL\n");
}
\end{lstlisting}
\risultato
\begin{lstlisting}[style=c, name=lista.c]
int lunghezza(struct Node *head)
{
   struct Node *p;
   int n = 0;

   for (p = head; p != NULL; p = p->next)
      n++;
   return n;
}
\end{lstlisting}
\end{affianco}

Qui \code{head} è un parametro, cioè una copia del puntatore del chiamante:
la funzione potrebbe anche scorrere la lista usando direttamente
\code{head}, senza cambiare la \code{head} di chi la chiama. Usare \code{p}
rende però il codice più chiaro, ed è l'abitudine giusta.

\section{Le operazioni sulle liste}

\subsection{Un modulo per le liste}

Le slide elencano le operazioni fondamentali: creazione della lista,
inserimento di un nodo, eliminazione di un nodo, ricerca di un elemento,
lunghezza della lista, unione di due liste. Le scriviamo come funzioni e le
raccogliamo in un \textbf{modulo}, come nel capitolo~\ref{chap:make}: i
prototipi in \code{lista.h}, le definizioni in \code{lista.c}. Un programma
che usa le liste include \code{lista.h} e si compila insieme a
\code{lista.c}.

\begin{lstlisting}[style=h, name=lista.h]
struct Node {
   int data;               /* il dato */
   struct Node *next;      /* il puntatore al nodo successivo */
};

struct Node *nuovoNodo(int data);
void stampaLista(struct Node *head);
int lunghezza(struct Node *head);
struct Node *inserisciTesta(struct Node *head, int data);
struct Node *inserisciCoda(struct Node *head, int data);
struct Node *inserisciPosizione(struct Node *head, int data, int pos);
struct Node *inserisciOrdinato(struct Node *head, int data);
struct Node *cerca(struct Node *head, int data);
struct Node *eliminaTesta(struct Node *head);
struct Node *eliminaCoda(struct Node *head);
struct Node *eliminaValore(struct Node *head, int data);
struct Node *eliminaPosizione(struct Node *head, int pos);
struct Node *unisci(struct Node *a, struct Node *b);
void liberaLista(struct Node *head);
\end{lstlisting}

\code{lista.c} comincia includendo \code{stdio.h}, \code{stdlib.h} (per
\code{malloc} e \code{free}) e \code{lista.h}; le sue funzioni sono mostrate
una alla volta nelle prossime sezioni. Sono state tutte provate anche nei
casi limite (lista vuota, posizioni che non esistono) con
\code{-fsanitize=address}, e con valgrind, che vedremo alla fine del
capitolo, per controllare che ogni nodo venga liberato.

\subsection{Restituire la nuova testa}

Molte operazioni possono cambiare il primo nodo: inserire in testa, eliminare
il primo nodo, inserire in una lista vuota. Ma \code{head} è una variabile
del chiamante, e la funzione ne riceve una \textbf{copia}
(capitolo~\ref{chap:funzioni}): assegnare un nuovo valore al parametro non
cambia la \code{head} del chiamante. Per questo le funzioni del modulo, come
quelle del progetto delle lezioni, \textbf{restituiscono la testa} della
lista dopo l'operazione, e il chiamante la riassegna:

\begin{lstlisting}[style=c]
head = inserisciTesta(head, 10);
\end{lstlisting}

È lo stesso schema di \code{e = update(e, 3)} del
capitolo~\ref{chap:strutture}. Un errore tipico è dimenticare
l'assegnazione, scrivendo soltanto \code{inserisciTesta(head, 10);}: il nodo
viene creato, ma \code{head} continua a puntare al vecchio primo nodo, e il
nuovo nodo è perso, con la sua memoria, che non potrà più essere liberata.

\subsection{Creare un nodo}

Ogni inserimento comincia creando un nodo nello heap:

\begin{affianco}[0.55][c]
\begin{lstlisting}[style=c, name=lista.c]
struct Node *nuovoNodo(int data)
{
   struct Node *nodo = malloc(sizeof(struct Node));

   if (nodo == NULL)
      return NULL;
   nodo->data = data;
   nodo->next = NULL;
   return nodo;
}
\end{lstlisting}
\risultato
\code{malloc} restituisce \code{NULL} se la memoria è finita: in quel caso
anche \code{nuovoNodo} restituisce \code{NULL}, e chi la chiama deve
controllarlo. Il campo \code{next} parte da \code{NULL}: il nodo appena
creato non è ancora collegato a niente.
\end{affianco}

\section{Inserimento}

Le slide prevedono quattro modi di inserire un nodo: in testa, in coda, in
una posizione specificata e in ordine (per dati che si possono ordinare,
come gli interi).

\subsection{In testa}

\begin{lstlisting}[style=c, name=lista.c]
struct Node *inserisciTesta(struct Node *head, int data)
{
   struct Node *nodo = nuovoNodo(data);

   if (nodo == NULL)
      return head;          /* niente memoria: la lista resta com'era */
   nodo->next = head;       /* 1. il nuovo nodo punta al vecchio primo */
   return nodo;             /* 2. e diventa il primo */
}
\end{lstlisting}

\begin{figure}[H]
  \centering
  \begin{tikzpicture}[font=\Lato\scriptsize]
    \variabile[minimum width=0.7cm]{h}{(0,0)}{head}{}
    \nodo[mlnuovo]{n1}{(1.0,1.05)}{10}
    \nodo{n2}{(3.1,0)}{20}
    \nodo{n3}{(5.0,0)}{30}
    \nodonull{n3}
    \collega{n2}{n3}
    \draw[mrifvecchio] (h.center) -- (n2.west);
    \draw[mrifnuovo] (n1-next.center) -- node[mpasso, pos=0.45, above right=0pt] {1} (n2.west);
    \draw[mrifnuovo] (h.center) -- node[mpasso, pos=0.5, above left=0pt] {2} (n1.west);
    \node[mpunto] at (n1-next.center) {};
    \node[mpunto] at (h.center) {};
    \node[mnota, anchor=south] at ([yshift=2pt]n1.north) {il nuovo nodo};
  \end{tikzpicture}
  \caption{Inserimento in testa: prima il nuovo nodo punta al vecchio primo
           nodo, poi \code{head} punta al nuovo nodo.}
  \label{fig:ins-testa}
\end{figure}

I passi sono due (figura~\ref{fig:ins-testa}): \wtpassonum{1} il nuovo nodo
punta al vecchio primo nodo, il cui indirizzo è in \code{head};
\wtpassonum{2} il nuovo nodo diventa il primo: la funzione ne restituisce
l'indirizzo, e il chiamante lo assegna a \code{head}. La funzione va bene
anche per la lista vuota: \code{nodo->next} riceve \code{NULL}, e il nuovo
nodo è l'unico della lista. Non c'è nessun ciclo: l'inserimento in testa
richiede lo stesso tempo, qualunque sia la lunghezza della lista.

\subsection{In coda}

Per inserire in coda serve l'ultimo nodo, cioè quello il cui \code{next}
vale \code{NULL}:

\begin{lstlisting}[style=c, name=lista.c]
struct Node *inserisciCoda(struct Node *head, int data)
{
   struct Node *nodo = nuovoNodo(data);
   struct Node *p;

   if (nodo == NULL)
      return head;
   if (head == NULL)        /* lista vuota: il nodo è l'unico */
      return nodo;
   for (p = head; p->next != NULL; p = p->next)
      ;                     /* si ferma sull'ultimo nodo */
   p->next = nodo;
   return head;
}
\end{lstlisting}

\begin{figure}[H]
  \centering
  \begin{tikzpicture}[font=\Lato\scriptsize]
    \variabile[minimum width=0.7cm]{h}{(0,0)}{head}{}
    \nodo{n1}{(0.8,0)}{10}
    \nodo{n2}{(2.7,0)}{20}
    \nodo{n3}{(4.6,0)}{30}
    \nodo[mlnuovo]{n4}{(6.9,0)}{40}
    \nodonull{n4}
    \riferimento{h}{-- (n1.west)}
    \collega{n1}{n2}
    \collega{n2}{n3}
    \draw[mrifnuovo] (n3-next.center) -- node[mpasso, above=1pt] {2} (n4.west);
    \node[mpunto] at (n3-next.center) {};
    \node[mvar, minimum width=0.6cm] (p) at ([yshift=-1.0cm]n3-dato.south) {};
    \node[mnome, anchor=east] (pn) at ([xshift=-2pt]p.west) {p};
    \riferimento{p}{-- (n3-dato.south)}
    \node[mpasso, anchor=east] at ([xshift=-3pt]pn.west) {1};
    \node[mnota, anchor=west] at ([xshift=4pt]p.east)
      {il ciclo si ferma sull'ultimo nodo: \code{p->next == NULL}};
    \node[mnota, anchor=south] at ([yshift=2pt]n4.north) {il nuovo nodo};
  \end{tikzpicture}
  \caption{Inserimento in coda: \code{p} si ferma sull'ultimo nodo, poi il
           suo \code{next} punta al nuovo nodo.}
  \label{fig:ins-coda}
\end{figure}

\begin{itemize}
  \item Il ciclo ha il corpo vuoto (il \code{;} da solo,
        capitolo~\ref{chap:flusso}): serve solo a far avanzare \code{p}. A
        differenza del ciclo di scorrimento, la condizione è
        \code{p->next != NULL}, quindi il ciclo si ferma \textbf{sull'ultimo
        nodo}, non dopo (figura~\ref{fig:ins-coda}). Poi \code{p->next = nodo}
        aggancia il nuovo nodo, il cui \code{next} è già \code{NULL}.
  \item La lista vuota va trattata a parte: non ha un ultimo nodo, e con
        \code{p} uguale a \code{NULL} la condizione \code{p->next != NULL}
        leggerebbe un campo di \code{NULL}, e il programma verrebbe fermato.
        In quel caso il nuovo nodo è tutta la lista.
  \item Per arrivare all'ultimo nodo bisogna attraversarli tutti:
        l'inserimento in coda richiede un tempo proporzionale alla lunghezza
        della lista.
\end{itemize}

\subsection{In una posizione}

Le posizioni si contano da 0, come gli indici degli array: inserire in
posizione 0 è inserire in testa. Per inserire in posizione \code{pos} bisogna
arrivare al nodo \textbf{precedente}, quello in posizione \code{pos - 1},
perché è il suo campo \code{next} che deve cambiare:

\begin{lstlisting}[style=c, name=lista.c]
struct Node *inserisciPosizione(struct Node *head, int data, int pos)
{
   struct Node *nodo, *prec;
   int i;

   if (pos == 0)
      return inserisciTesta(head, data);
   prec = head;             /* cerca il nodo in posizione pos - 1 */
   for (i = 0; prec != NULL && i < pos - 1; i++)
      prec = prec->next;
   if (prec == NULL)        /* la lista è troppo corta */
      return head;
   nodo = nuovoNodo(data);
   if (nodo == NULL)
      return head;
   nodo->next = prec->next; /* 1. il nuovo nodo punta al successivo */
   prec->next = nodo;       /* 2. il precedente punta al nuovo nodo */
   return head;
}
\end{lstlisting}

\begin{figure}[H]
  \centering
  \begin{tikzpicture}[font=\Lato\scriptsize]
    \variabile[minimum width=0.7cm]{h}{(0,0)}{head}{}
    \nodo{n1}{(0.8,0)}{10}
    \nodo{n2}{(2.7,0)}{20}
    \nodo{n4}{(6.0,0)}{40}
    \nodonull{n4}
    \nodo[mlnuovo]{n3}{(4.3,-1.15)}{30}
    \riferimento{h}{-- (n1.west)}
    \collega{n1}{n2}
    \draw[mrifvecchio] (n2-next.center) -- (n4.west);
    \draw[mrifnuovo] (n3-next.center) -- node[mpasso, pos=0.45, right=3pt] {1} (n4.west);
    \draw[mrifnuovo] (n2-next.center) -- node[mpasso, pos=0.5, right=2pt] {2} (n3.west);
    \node[mpunto] at (n2-next.center) {};
    \node[mpunto] at (n3-next.center) {};
    \node[mvar, minimum width=0.6cm] (pr) at ([yshift=1.0cm]n2-dato.north) {};
    \node[mnome, anchor=east] at ([xshift=-2pt]pr.west) {prec};
    \riferimento{pr}{-- (n2-dato.north)}
    \node[mindir, anchor=north] at ([yshift=-1pt]n1-dato.south) {posizione 0};
    \node[mindir, anchor=north] at ([yshift=-1pt]n2-dato.south) {posizione 1};
    \node[mnota, anchor=north] at ([yshift=-2pt]n3.south) {il nuovo nodo, in posizione 2};
  \end{tikzpicture}
  \caption{Inserimento di 30 in posizione 2: \code{prec} si ferma sul nodo
           in posizione 1; poi il nuovo nodo punta al successivo di
           \code{prec}, e \code{prec} punta al nuovo nodo.}
  \label{fig:ins-pos}
\end{figure}

Il ciclo \code{for} fa avanzare \code{prec} di \code{pos - 1} nodi a partire
dal primo; se la lista finisce prima (\code{prec} diventa \code{NULL}), la
posizione non esiste, e la lista resta com'è. Poi vengono i due passi della
figura~\ref{fig:ins-pos}: \wtpassonum{1} il nuovo nodo punta al successivo
di \code{prec}; \wtpassonum{2} \code{prec} punta al nuovo nodo. La testa non
cambia, e la funzione restituisce \code{head} così com'era.

\warning{L'ordine dei due passi conta}{
  Se si scambiano i passi, scrivendo prima \code{prec->next = nodo} e poi
  \code{nodo->next = prec->next}, il secondo assegnamento copia in
  \code{nodo->next} l'indirizzo di \code{nodo} stesso: il nuovo nodo punta
  a se stesso, e l'indirizzo del resto della lista, che stava solo in
  \code{prec->next}, è perso per sempre. Uno scorrimento della lista non
  terminerebbe più. La regola: prima si collega il nuovo nodo al resto
  della lista, poi lo si aggancia al precedente.}

\subsection{In ordine}

Se la lista è ordinata (qui in ordine crescente), un nuovo dato va inserito
al suo posto, in modo che la lista resti ordinata. È un inserimento in
posizione, in cui la posizione si trova confrontando i dati:

\begin{lstlisting}[style=c, name=lista.c]
struct Node *inserisciOrdinato(struct Node *head, int data)
{
   struct Node *nodo, *prec;

   if (head == NULL || data <= head->data)
      return inserisciTesta(head, data);
   prec = head;             /* avanza finché il successivo è minore */
   while (prec->next != NULL && prec->next->data < data)
      prec = prec->next;
   nodo = nuovoNodo(data);
   if (nodo == NULL)
      return head;
   nodo->next = prec->next;
   prec->next = nodo;
   return head;
}
\end{lstlisting}

\begin{itemize}
  \item Se la lista è vuota, o il nuovo dato non supera il primo, il nodo
        va in testa. (Anche qui conta l'ordine delle condizioni: se
        \code{head} è \code{NULL}, la valutazione in cortocircuito,
        capitolo~\ref{chap:flusso}, evita di leggere \code{head->data}.)
  \item Altrimenti \code{prec} avanza finché il nodo \textbf{successivo} ha
        un dato minore di quello da inserire; quando si ferma, il nuovo nodo
        va subito dopo \code{prec}, con gli stessi due passi
        dell'inserimento in posizione.
  \item Se \code{prec} arriva all'ultimo nodo, \code{prec->next} vale
        \code{NULL}: la prima condizione del \code{while} è falsa, la
        seconda non viene valutata (non si legge il campo \code{data} di
        \code{NULL}), e il nuovo nodo va in coda.
\end{itemize}

Seguiamo a mano l'inserimento di 25 nella lista 10, 20, 30, 40. Il primo
dato, 10, è minore di 25, quindi non si inserisce in testa, e \code{prec}
parte dal primo nodo:

\begin{center}
\renewcommand{\arraystretch}{1.15}
\begin{tabular}{@{}cccl@{}}
  \toprule
  \textbf{\code{prec}} & \textbf{\code{prec->next->data}} & \textbf{minore di 25?} & \textbf{che cosa succede} \\
  \midrule
  10 & 20 & sì & \code{prec} avanza \\
  20 & 30 & no & il ciclo si ferma: 25 va fra 20 e 30 \\
  \bottomrule
\end{tabular}
\end{center}

\subsection{Una prova}

Il programma seguente usa le quattro funzioni, partendo da una lista vuota:

\begin{affianco}[0.6][c]
\begin{lstlisting}[style=c, name=inserimenti.c]
#include <stdio.h>
#include "lista.h"

int main(void)
{
   struct Node *head = NULL;

   head = inserisciTesta(head, 20);
   head = inserisciTesta(head, 10);
   printf("in testa 20 e 10:    ");
   stampaLista(head);
   head = inserisciCoda(head, 40);
   printf("in coda 40:          ");
   stampaLista(head);
   head = inserisciPosizione(head, 30, 2);
   printf("30 in posizione 2:   ");
   stampaLista(head);
   head = inserisciOrdinato(head, 25);
   head = inserisciOrdinato(head, 5);
   head = inserisciOrdinato(head, 50);
   printf("in ordine 25, 5, 50: ");
   stampaLista(head);
   printf("lunghezza: %d\n", lunghezza(head));
   liberaLista(head);
   return 0;
}
\end{lstlisting}
\risultato
\begin{itemize}
  \item Inserendo in testa prima 20 e poi 10, la lista è 10, 20: ogni nodo
        inserito in testa va davanti a tutti gli altri, quindi l'ordine
        della lista è l'opposto di quello degli inserimenti.
  \item 30 va in posizione 2, fra 20 e 40.
  \item Con \code{inserisciOrdinato} il 25 va fra 20 e 30, il 5 in testa e
        il 50 in coda: la lista resta ordinata.
  \item \code{liberaLista}, che vedremo alla fine, restituisce la memoria
        di tutti i nodi.
\end{itemize}
\end{affianco}

\begin{console}
$ gcc -Wall -o inserimenti inserimenti.c lista.c
$ ./inserimenti
in testa 20 e 10:    10 -> 20 -> NULL
in coda 40:          10 -> 20 -> 40 -> NULL
30 in posizione 2:   10 -> 20 -> 30 -> 40 -> NULL
in ordine 25, 5, 50: 5 -> 10 -> 20 -> 25 -> 30 -> 40 -> 50 -> NULL
lunghezza: 7
\end{console}

\section{Ricerca}

Cercare un elemento vuol dire trovare il nodo che contiene un certo dato. La
funzione scorre la lista e restituisce l'\textbf{indirizzo} del primo nodo
con quel dato, oppure \code{NULL} se non c'è:

\begin{affianco}[0.55][c]
\begin{lstlisting}[style=c, name=lista.c]
struct Node *cerca(struct Node *head, int data)
{
   struct Node *p;

   for (p = head; p != NULL; p = p->next)
      if (p->data == data)
         return p;          /* trovato: l'indirizzo del nodo */
   return NULL;             /* non trovato */
}
\end{lstlisting}
\risultato
Restituire l'indirizzo, e non un semplice sì o no, permette al chiamante di
usare il nodo trovato: leggerne o cambiarne i dati, guardare il successivo.
\end{affianco}

La ricerca è necessariamente \textbf{sequenziale}: anche se la lista è
ordinata, non c'è modo di saltare a metà lista, e nel caso peggiore (dato
assente, o nell'ultimo nodo) bisogna attraversarla tutta. Una prova della
ricerca è nella sezione seguente, insieme a quelle delle eliminazioni.

\section{Eliminazione}

Eliminare un nodo richiede due cose: \textbf{scollegarlo}, cioè fare in modo
che nessun nodo della lista punti più a lui, e \textbf{liberarne} la memoria
con \code{free}. Le slide elencano quattro casi: in testa, in coda, per
valore (il primo nodo che ha una certa proprietà, qui un certo dato) e data
la posizione.

\subsection{In testa}

\begin{lstlisting}[style=c, name=lista.c]
struct Node *eliminaTesta(struct Node *head)
{
   struct Node *secondo;

   if (head == NULL)
      return NULL;
   secondo = head->next;    /* 1. ricorda il secondo nodo */
   free(head);              /* 2. libera il primo */
   return secondo;          /* 3. il secondo diventa il primo */
}
\end{lstlisting}

\begin{figure}[H]
  \centering
  \begin{tikzpicture}[font=\Lato\scriptsize]
    \variabile[minimum width=0.7cm]{h}{(0,0)}{head}{}
    \nodo[mlvia]{n1}{(0.9,0)}{10}
    \nodo{n2}{(2.9,0)}{20}
    \nodo{n3}{(4.8,0)}{30}
    \nodonull{n3}
    \collega[mrifvecchio]{n1}{n2}
    \collega{n2}{n3}
    \draw[mrifvecchio] (h.center) -- (n1.west);
    \draw[mrifnuovo] (h.center) to[bend right=40] node[mpasso, pos=0.5, below=1pt] {3} (n2.south west);
    \node[mpunto] at (h.center) {};
    \node[mvar, minimum width=0.6cm] (s) at ([yshift=1.0cm]n2-dato.north) {};
    \node[mnome, anchor=east] at ([xshift=-2pt]s.west) {secondo};
    \riferimento{s}{-- (n2-dato.north)}
    \node[mpasso, anchor=west] at ([xshift=3pt]s.east) {1};
    \node[mpasso] (q2) at ([yshift=0.45cm]n1-dato.north) {2};
    \node[mnota, anchor=east] at ([xshift=-1pt]q2.west) {\code{free(head)}};
  \end{tikzpicture}
  \caption{Eliminazione in testa: si ricorda il secondo nodo, si libera il
           primo, e il secondo diventa la nuova testa.}
  \label{fig:el-testa}
\end{figure}

Il passo \wtpassonum{1} è indispensabile: dopo \code{free(head)} il nodo non
esiste più, e il suo campo \code{next} non si può più leggere. L'indirizzo
del secondo nodo va quindi salvato \textbf{prima}. Poi il chiamante assegna
il risultato a \code{head} (passo \wtpassonum{3}, figura~\ref{fig:el-testa}).

\subsection{In coda}

\begin{lstlisting}[style=c, name=lista.c]
struct Node *eliminaCoda(struct Node *head)
{
   struct Node *p;

   if (head == NULL)
      return NULL;
   if (head->next == NULL) {   /* un solo nodo */
      free(head);
      return NULL;
   }
   for (p = head; p->next->next != NULL; p = p->next)
      ;                     /* si ferma sul penultimo nodo */
   free(p->next);
   p->next = NULL;
   return head;
}
\end{lstlisting}

\begin{figure}[H]
  \centering
  \begin{tikzpicture}[font=\Lato\scriptsize]
    \variabile[minimum width=0.7cm]{h}{(0,0)}{head}{}
    \nodo{n1}{(0.8,0)}{10}
    \nodo{n2}{(2.7,0)}{20}
    \nodo[mlvia]{n3}{(5.0,0)}{30}
    \nodonull{n3}
    \riferimento{h}{-- (n1.west)}
    \collega{n1}{n2}
    \draw[mrifvecchio] (n2-next.east) -- (n3.west);
    \draw[line width=0.8pt, draw=mok] ([shift={(1.5pt,1.5pt)}]n2-next.south west)
      -- ([shift={(-1.5pt,-1.5pt)}]n2-next.north east);
    \node[mvar, minimum width=0.6cm] (p) at ([yshift=-1.0cm]n2-dato.south) {};
    \node[mnome, anchor=east] (pn) at ([xshift=-2pt]p.west) {p};
    \riferimento{p}{-- (n2-dato.south)}
    \node[mpasso, anchor=east] at ([xshift=-3pt]pn.west) {1};
    \node[mnota, anchor=west] at ([xshift=4pt]p.east)
      {si ferma sul penultimo nodo: \code{p->next->next == NULL}};
    \node[mpasso] (q2) at ([yshift=0.45cm]n3-dato.north) {2};
    \node[mnota, anchor=west] at ([xshift=1pt]q2.east) {\code{free(p->next)}};
    \node[mpasso] (q3) at ([yshift=0.45cm]n2-next.north) {3};
    \node[mnota, anchor=east] at ([xshift=-1pt]q3.west) {\code{p->next = NULL}};
  \end{tikzpicture}
  \caption{Eliminazione in coda: \code{p} si ferma sul penultimo nodo, che
           dopo il \code{free} dell'ultimo diventa l'ultimo.}
  \label{fig:el-coda}
\end{figure}

Per eliminare l'ultimo nodo bisogna cambiare il \code{next} del
\textbf{penultimo}, che deve diventare \code{NULL}. Il ciclo si ferma quando
il successivo del successivo di \code{p} non esiste:
\code{p->next->next == NULL} vuol dire che \code{p->next} è l'ultimo nodo,
e quindi \code{p} è il penultimo (figura~\ref{fig:el-coda}). Una lista di un
solo nodo non ha un penultimo, e va trattata a parte: dopo l'eliminazione è
vuota, e la funzione restituisce \code{NULL}.

\subsection{Per valore}

\begin{lstlisting}[style=c, name=lista.c]
struct Node *eliminaValore(struct Node *head, int data)
{
   struct Node *prec, *p;

   if (head == NULL)
      return NULL;
   if (head->data == data)
      return eliminaTesta(head);
   prec = head;             /* avanza finché il successivo non è quello */
   while (prec->next != NULL && prec->next->data != data)
      prec = prec->next;
   if (prec->next != NULL) {   /* trovato: è prec->next */
      p = prec->next;
      prec->next = p->next;    /* 1. il precedente scavalca il nodo */
      free(p);                 /* 2. il nodo viene liberato */
   }
   return head;
}
\end{lstlisting}

\begin{figure}[H]
  \centering
  \begin{tikzpicture}[font=\Lato\scriptsize]
    \variabile[minimum width=0.7cm]{h}{(0,0)}{head}{}
    \nodo{n1}{(0.8,0)}{20}
    \nodo{n2}{(2.7,0)}{30}
    \nodo[mlvia]{n3}{(4.6,0)}{40}
    \nodo{n4}{(6.5,0)}{50}
    \nodonull{n4}
    \riferimento{h}{-- (n1.west)}
    \collega{n1}{n2}
    \collega[mrifvecchio]{n2}{n3}
    \collega[mrifvecchio]{n3}{n4}
    \draw[mrifnuovo] (n2-next.center) to[bend right=45]
      node[mpasso, pos=0.5, below=1pt] {1} (n4.south west);
    \node[mvar, minimum width=0.6cm] (pr) at ([yshift=1.0cm]n2-dato.north) {};
    \node[mnome, anchor=east] at ([xshift=-2pt]pr.west) {prec};
    \riferimento{pr}{-- (n2-dato.north)}
    \node[mvar, minimum width=0.6cm] (pp) at ([yshift=1.0cm]n3-dato.north) {};
    \node[mnome, anchor=east] at ([xshift=-2pt]pp.west) {p};
    \riferimento{pp}{-- (n3-dato.north)}
    \node[mpasso, anchor=west] (q2) at ([xshift=4pt]pp.east) {2};
    \node[mnota, anchor=west] at ([xshift=1pt]q2.east) {\code{free(p)}};
  \end{tikzpicture}
  \caption{Eliminazione del 40: \code{prec} scavalca il nodo da eliminare,
           puntando al successivo, e poi il nodo viene liberato.}
  \label{fig:el-valore}
\end{figure}

Come per l'inserimento, serve il nodo \textbf{precedente} a quello da
eliminare, perché è il suo \code{next} che cambia. Il ciclo avanza finché il
successivo di \code{prec} non contiene il dato cercato; se lo trova,
\wtpassonum{1} \code{prec} scavalca il nodo, puntando a quello dopo, e
\wtpassonum{2} il nodo viene liberato (figura~\ref{fig:el-valore}). Anche
qui l'ordine conta: dopo \code{free(p)} non si potrebbe più leggere
\code{p->next}. Se il dato non c'è, la lista resta com'è; se è nel primo
nodo, è un'eliminazione in testa.

\subsection{Data la posizione}

\begin{lstlisting}[style=c, name=lista.c]
struct Node *eliminaPosizione(struct Node *head, int pos)
{
   struct Node *prec, *p;
   int i;

   if (head == NULL)
      return NULL;
   if (pos == 0)
      return eliminaTesta(head);
   prec = head;             /* cerca il nodo in posizione pos - 1 */
   for (i = 0; prec->next != NULL && i < pos - 1; i++)
      prec = prec->next;
   if (prec->next == NULL)  /* non c'è un nodo in posizione pos */
      return head;
   p = prec->next;
   prec->next = p->next;
   free(p);
   return head;
}
\end{lstlisting}

È la stessa operazione, ma \code{prec} si ferma sul nodo in posizione
\code{pos - 1} contando i nodi, invece di confrontare i dati. Se il nodo in
posizione \code{pos} non esiste, la lista resta com'è.

\subsection{Una prova}

\begin{affianco}[0.6][c]
\begin{lstlisting}[style=c, name=eliminazioni.c]
#include <stdio.h>
#include "lista.h"

int main(void)
{
   struct Node *head = NULL, *p;
   int i;

   for (i = 1; i <= 7; i++)
      head = inserisciCoda(head, i * 10);
   printf("all'inizio:          ");
   stampaLista(head);
   p = cerca(head, 40);
   if (p != NULL)
      printf("cerca(40): trovato, il successivo vale %d\n",
             p->next->data);
   if (cerca(head, 45) == NULL)
      printf("cerca(45): non trovato\n");
   head = eliminaTesta(head);
   printf("elimina in testa:    ");
   stampaLista(head);
   head = eliminaCoda(head);
   printf("elimina in coda:     ");
   stampaLista(head);
   head = eliminaValore(head, 40);
   printf("elimina il 40:       ");
   stampaLista(head);
   head = eliminaPosizione(head, 1);
   printf("elimina posizione 1: ");
   stampaLista(head);
   liberaLista(head);
   return 0;
}
\end{lstlisting}
\risultato
\begin{itemize}
  \item \code{cerca(head, 40)} restituisce l'indirizzo del nodo con il 40,
        e attraverso di esso il programma legge il dato del nodo
        successivo, 50.
  \item Le eliminazioni tolgono, nell'ordine, il primo nodo (10), l'ultimo
        (70), il nodo con il 40 e quello in posizione 1, che a quel punto è
        il 30.
\end{itemize}
\end{affianco}

\begin{console}
$ gcc -Wall -o eliminazioni eliminazioni.c lista.c
$ ./eliminazioni
all'inizio:          10 -> 20 -> 30 -> 40 -> 50 -> 60 -> 70 -> NULL
cerca(40): trovato, il successivo vale 50
cerca(45): non trovato
elimina in testa:    20 -> 30 -> 40 -> 50 -> 60 -> 70 -> NULL
elimina in coda:     20 -> 30 -> 40 -> 50 -> 60 -> NULL
elimina il 40:       20 -> 30 -> 50 -> 60 -> NULL
elimina posizione 1: 20 -> 50 -> 60 -> NULL
\end{console}

\section{Unione di due liste}

\subsection{Agganciare due liste}

L'unione (concatenazione) di due liste \code{a} e \code{b} si ottiene
agganciando il primo nodo di \code{b} all'ultimo nodo di \code{a}: non si
copia nessun nodo, si cambia un solo puntatore (figura~\ref{fig:unione}). Se
\code{a} è vuota, l'unione è semplicemente \code{b}.

\begin{affianco}[0.55][c]
\begin{lstlisting}[style=c, name=lista.c]
struct Node *unisci(struct Node *a, struct Node *b)
{
   struct Node *p;

   if (a == NULL)
      return b;
   for (p = a; p->next != NULL; p = p->next)
      ;                     /* l'ultimo nodo di a */
   p->next = b;             /* punta al primo nodo di b */
   return a;
}
\end{lstlisting}
\risultato
\begin{lstlisting}[style=c, name=unione.c]
#include <stdio.h>
#include "lista.h"

int main(void)
{
   struct Node *a = NULL, *b = NULL;

   a = inserisciCoda(a, 1);
   a = inserisciCoda(a, 2);
   a = inserisciCoda(a, 3);
   b = inserisciCoda(b, 7);
   b = inserisciCoda(b, 8);
   printf("a: ");
   stampaLista(a);
   printf("b: ");
   stampaLista(b);
   a = unisci(a, b);
   printf("a dopo unisci(a, b): ");
   stampaLista(a);
   printf("lunghezza: %d\n", lunghezza(a));
   liberaLista(a);    /* anche i nodi di b */
   return 0;
}
\end{lstlisting}
\end{affianco}

\begin{figure}[H]
  \centering
  \begin{tikzpicture}[font=\Lato\scriptsize]
    \variabile[minimum width=0.7cm]{ha}{(0,0)}{a}{}
    \nodo{a1}{(0.8,0)}{1}
    \nodo{a2}{(2.7,0)}{2}
    \nodo{a3}{(4.6,0)}{3}
    \riferimento{ha}{-- (a1.west)}
    \collega{a1}{a2}
    \collega{a2}{a3}
    \variabile[minimum width=0.7cm]{hb}{(5.75,-1.3)}{b}{}
    \nodo{b1}{(6.55,-1.3)}{7}
    \nodo{b2}{(8.45,-1.3)}{8}
    \nodonull{b2}
    \riferimento{hb}{-- (b1.west)}
    \collega{b1}{b2}
    \draw[mrifnuovo] (a3-next.center) -| node[mnota, pos=0.75, right=2pt] {\code{p->next = b}} (b1-dato.north);
    \node[mpunto] at (a3-next.center) {};
  \end{tikzpicture}
  \caption{\code{unisci(a, b)}: il \code{next} dell'ultimo nodo di \code{a},
           che valeva \code{NULL}, punta al primo nodo di \code{b}.}
  \label{fig:unione}
\end{figure}

\begin{console}
$ gcc -Wall -o unione unione.c lista.c
$ ./unione
a: 1 -> 2 -> 3 -> NULL
b: 7 -> 8 -> NULL
a dopo unisci(a, b): 1 -> 2 -> 3 -> 7 -> 8 -> NULL
lunghezza: 5
\end{console}

Dopo l'unione i nodi di \code{b} fanno parte di \code{a}, e \code{b} punta
ancora al nodo 7, che ora sta in mezzo alla lista \code{a}. Per questo il
programma libera solo \code{a}: liberare anche \code{b} vorrebbe dire
liberare due volte gli stessi nodi (\textit{double free}), un errore grave
quanto usare un nodo già liberato.

\section{Liberare la lista}

\subsection{Un nodo alla volta}

Quando una lista non serve più, ogni suo nodo va restituito con \code{free}:
è un'eliminazione in testa (figura~\ref{fig:el-testa}) ripetuta finché la
lista non è vuota.

\begin{affianco}[0.55][c]
\begin{lstlisting}[style=c, name=lista.c]
void liberaLista(struct Node *head)
{
   struct Node *succ;

   while (head != NULL) {
      succ = head->next;    /* salva il successivo */
      free(head);
      head = succ;
   }
}
\end{lstlisting}
\risultato
Il punto delicato è lo stesso dell'eliminazione: dopo \code{free(head)} il
nodo non esiste più, e il suo campo \code{next} non si può leggere.
L'indirizzo del successivo si salva quindi \textbf{prima}, in \code{succ}.
\end{affianco}

\subsection{L'errore tipico}

Sembra più semplice scrivere il ciclo di scorrimento con un \code{free} nel
corpo, come in \code{liberaSbagliata}:

\begin{lstlisting}[style=c, name=libera_male.c]
#include <stdlib.h>
#include "lista.h"

void liberaSbagliata(struct Node *head)
{
   struct Node *p;

   for (p = head; p != NULL; p = p->next)
      free(p);             /* poi p = p->next legge un nodo già liberato */
}

int main(void)
{
   struct Node *head = NULL;

   head = inserisciCoda(head, 10);
   head = inserisciCoda(head, 20);
   liberaSbagliata(head);
   return 0;
}
\end{lstlisting}

Ma \code{p = p->next} viene eseguito \textbf{dopo} \code{free(p)}: legge il
campo \code{next} di un nodo che non esiste più. gcc 13 se ne accorge già in
compilazione, con \code{-Wall}:

\begin{console}
$ gcc -Wall -o libera_male libera_male.c lista.c
libera_male.c: In function (*@‘@*)liberaSbagliata(*@’@*):
libera_male.c:8:32: warning: pointer (*@‘@*)p(*@’@*) used after (*@‘@*)free(*@’@*) [-Wuse-after-free]
    8 |    for (p = head; p != NULL; p = p->next)
      |                              ~~^~~~~~~~~
libera_male.c:9:7: note: call to (*@‘@*)free(*@’@*) here
    9 |       free(p);             /* poi p = p->next legge un nodo già liberato */
      |       ^~~~~~~
$ ./libera_male
Segmentation fault
\end{console}

Che cosa succede in esecuzione dipende da che cosa c'è nella memoria
liberata: nella prova, \code{free} aveva già usato quei byte per le proprie
informazioni interne, il ciclo ha letto come \code{next} un indirizzo senza
senso, e il programma è stato fermato. Altre volte un errore del genere
passa inosservato, e per questo è insidioso. Con \code{-fsanitize=address}
(capitolo~\ref{chap:puntatori}) il rapporto è preciso (qui abbreviato, con
un segnaposto al posto del percorso della cartella):

\begin{console}
$ gcc -Wall -g -fsanitize=address -o libera_male libera_male.c lista.c
(*@\textit{\dots{} lo stesso avviso di prima \dots}@*)
$ ./libera_male
=================================================================
==4236==ERROR: AddressSanitizer: heap-use-after-free on address 0x502000000018 at pc 0x5645579ac2f4 bp 0x7ffe8d8b8390 sp 0x7ffe8d8b8380
READ of size 8 at 0x502000000018 thread T0
    #0 0x5645579ac2f3 in liberaSbagliata (*@\textit{[percorso del file]}@*)/libera_male.c:8
    #1 0x5645579ac354 in main (*@\textit{[percorso del file]}@*)/libera_male.c:18
(*@\textit{\dots}@*)
0x502000000018 is located 8 bytes inside of 16-byte region [0x502000000010,0x502000000020)
freed by thread T0 here:
    #0 0x7f8074efc4d8 in free ../../../../src/libsanitizer/asan/asan_malloc_linux.cpp:52
    #1 0x5645579ac2ce in liberaSbagliata (*@\textit{[percorso del file]}@*)/libera_male.c:9
(*@\textit{\dots}@*)
\end{console}

Il rapporto si legge così: alla riga 8 si è letto (\textit{READ}) un valore
di 8 byte, un puntatore, a 8 byte dall'inizio di un blocco di 16 byte, cioè
proprio il campo \code{next} di un nodo; e quel blocco era stato liberato
alla riga 9.

\section{La ricorsione}

\subsection{Una funzione che chiama se stessa}

Una funzione è \textbf{ricorsiva} quando, per risolvere un problema, chiama
se stessa su un problema più piccolo dello stesso tipo. Con le liste viene
naturale, perché una lista ha una struttura ricorsiva: una lista

\begin{itemize}
  \item o è \textbf{vuota} (\code{NULL});
  \item oppure è un \textbf{nodo seguito da una lista}: quella che comincia
        da \code{head->next}, che ha un nodo in meno.
\end{itemize}

Ne segue una definizione ricorsiva della lunghezza: la lunghezza della lista
vuota è 0; la lunghezza di un nodo seguito da una lista è 1 più la lunghezza
di quella lista. In C si scrive quasi parola per parola:

\begin{lstlisting}[style=c, name=ricorsione.c]
int lunghezzaRic(struct Node *head)
{
   if (head == NULL)               /* caso base: la lista vuota */
      return 0;
   return 1 + lunghezzaRic(head->next);
}
\end{lstlisting}

Ogni funzione ricorsiva ha due ingredienti:

\begin{itemize}
  \item il \textbf{caso base}, che si risolve senza altre chiamate (qui la
        lista vuota, la cui lunghezza è 0);
  \item il \textbf{passo ricorsivo}, che riduce il problema e chiama la
        funzione sul problema ridotto (qui il resto della lista, che ha un
        nodo in meno).
\end{itemize}

Poiché a ogni chiamata la lista si accorcia di un nodo, prima o poi si arriva
alla lista vuota, e le chiamate si fermano. Senza il caso base, arrivata in
fondo alla lista la funzione leggerebbe \code{head->next} con \code{head}
uguale a \code{NULL}, e il programma verrebbe fermato.

\subsection{Una traccia delle chiamate}

Che una funzione chiami se stessa non ha niente di speciale: ogni chiamata è
indipendente dalle altre, e ha la sua copia dei parametri e delle variabili
locali sulla pila delle chiamate (capitolo~\ref{chap:funzioni}). La versione
seguente di \code{lunghezzaRic} racconta che cosa succede, con un rientro
che cresce a ogni chiamata aperta: la variabile globale \code{livello} conta
le chiamate aperte, e \code{rientra} stampa quattro spazi per ognuna.

\begin{lstlisting}[style=c, name=traccia.c]
#include <stdio.h>
#include "lista.h"

int livello = 0;                   /* quante chiamate sono aperte */

void rientra(void)
{
   int i;

   for (i = 0; i < livello; i++)
      printf("    ");
}

int lunghezzaRic(struct Node *head)
{
   int n;

   rientra();
   if (head == NULL) {
      printf("lunghezzaRic(NULL): caso base, restituisce 0\n");
      return 0;
   }
   printf("lunghezzaRic(lista da %d): chiama lunghezzaRic sul resto\n",
          head->data);
   livello++;
   n = 1 + lunghezzaRic(head->next);
   livello--;
   rientra();
   printf("lunghezzaRic(lista da %d): restituisce %d\n", head->data, n);
   return n;
}

int main(void)
{
   struct Node *head = NULL;

   head = inserisciCoda(head, 10);
   head = inserisciCoda(head, 20);
   head = inserisciCoda(head, 30);
   printf("risultato: %d\n", lunghezzaRic(head));
   liberaLista(head);
   return 0;
}
\end{lstlisting}

\begin{console}
$ gcc -Wall -o traccia traccia.c lista.c
$ ./traccia
lunghezzaRic(lista da 10): chiama lunghezzaRic sul resto
    lunghezzaRic(lista da 20): chiama lunghezzaRic sul resto
        lunghezzaRic(lista da 30): chiama lunghezzaRic sul resto
            lunghezzaRic(NULL): caso base, restituisce 0
        lunghezzaRic(lista da 30): restituisce 1
    lunghezzaRic(lista da 20): restituisce 2
lunghezzaRic(lista da 10): restituisce 3
risultato: 3
\end{console}

\subsection{Che cosa succede}

\begin{itemize}
  \item La prima chiamata riceve la lista intera, che comincia da 10. Non è
        vuota, quindi deve calcolare il valore di
        \code{1 + lunghezzaRic(head->next)}, e per farlo chiama la funzione
        sulla lista che comincia da 20. Il
        calcolo della prima chiamata resta in sospeso: la chiamata è
        \textbf{aperta}, e la sua copia di \code{head} resta sulla pila.
  \item Lo stesso succede per la lista da 20 e per quella da 30: alla fine
        quattro chiamate sono aperte insieme, ognuna con il suo \code{head}
        (figura~\ref{fig:pila-ric}).
  \item La quarta chiamata riceve \code{NULL}: è il caso base, e restituisce
        0 senza chiamare nessuno.
  \item Ora le chiamate si chiudono a ritroso: quella sulla lista da 30
        completa \code{1 + 0} e restituisce 1; quella da 20 restituisce 2; la
        prima restituisce 3 a \code{main}.
\end{itemize}

\begin{figure}[H]
  \centering
  \begin{tikzpicture}[font=\Lato\scriptsize]
    \nodo{n1}{(0,0)}{10}
    \nodo{n2}{(2.4,0)}{20}
    \nodo{n3}{(4.8,0)}{30}
    \nodonull{n3}
    \collega{n1}{n2}
    \collega{n2}{n3}
    \foreach \k/\ix/\iy in {1/0.475/-1.45, 2/2.875/-2.65, 3/5.275/-3.85} {
      \node[mvar, minimum width=0.6cm] (h\k) at (\ix,\iy) {};
      \node[mnome, anchor=west] (h\k-n) at ([xshift=2pt]h\k.east) {head};
      \node[mscopolab, anchor=north west] (h\k-l) at ([xshift=-2pt, yshift=-3pt]h\k.south west) {chiamata \k};
      \begin{pgfonlayer}{liv2}
        \node[mscopo2, fit=(h\k)(h\k-n)(h\k-l)] (f\k) {};
      \end{pgfonlayer}
      \riferimento{h\k}{-- (n\k-dato.south)}
    }
    \node[mvar, minimum width=0.9cm, font=\FiraCodeNL\tiny] (h4) at (7.675,-5.05) {NULL};
    \node[mnome, anchor=west] (h4-n) at ([xshift=2pt]h4.east) {head};
    \node[mscopolab, anchor=north west] (h4-l) at ([xshift=-2pt, yshift=-3pt]h4.south west) {chiamata 4: caso base};
    \begin{pgfonlayer}{liv2}
      \node[mscopo2, fit=(h4)(h4-n)(h4-l)] (f4) {};
    \end{pgfonlayer}
    % i valori restituiti
    \draw[mcerca] (f4.west) -| node[mnota, pos=0.25, below] {restituisce 0} (f3.south);
    \draw[mcerca] (f3.west) -| node[mnota, pos=0.25, below] {restituisce 1} (f2.south);
    \draw[mcerca] (f2.west) -| node[mnota, pos=0.25, below] {restituisce 2} (f1.south);
    \draw[mcerca] (f1.west) -- ++(-0.8,0) node[mnota, anchor=east] {3 a \code{main}};
  \end{tikzpicture}
  \caption{Le quattro chiamate di \code{lunghezzaRic} aperte insieme, ognuna
           con la sua copia di \code{head}. Le frecce tratteggiate sono i
           valori restituiti, dall'ultima chiamata alla prima.}
  \label{fig:pila-ric}
\end{figure}

\subsection{Altri esempi}

Il programma \code{ricorsione.c} contiene, oltre a \code{lunghezzaRic}, tre
funzioni ricorsive sulle liste, e le usa su una lista di quattro nodi:

\begin{lstlisting}[style=c, name=ricorsione.c]
#include <stdio.h>
#include <stdlib.h>
#include "lista.h"

int lunghezzaRic(struct Node *head)
{
   if (head == NULL)               /* caso base: la lista vuota */
      return 0;
   return 1 + lunghezzaRic(head->next);
}

void stampaRovescia(struct Node *head)
{
   if (head == NULL)
      return;
   stampaRovescia(head->next);     /* prima il resto della lista */
   printf("%d ", head->data);      /* poi questo nodo */
}

struct Node *inserisciCodaRic(struct Node *head, struct Node *nodo)
{
   if (head == NULL)               /* in fondo: qui va il nodo */
      return nodo;
   head->next = inserisciCodaRic(head->next, nodo);
   return head;
}

void liberaListaRic(struct Node *head)
{
   if (head != NULL) {
      liberaListaRic(head->next);  /* prima il resto della lista */
      free(head);                  /* poi questo nodo */
   }
}

int main(void)
{
   struct Node *head = NULL;
   int i;

   for (i = 1; i <= 4; i++)
      head = inserisciCodaRic(head, nuovoNodo(i * 10));
   stampaLista(head);
   printf("lunghezza: %d\n", lunghezzaRic(head));
   printf("al contrario: ");
   stampaRovescia(head);
   printf("\n");
   liberaListaRic(head);
   return 0;
}
\end{lstlisting}

\begin{console}
$ gcc -Wall -o ricorsione ricorsione.c lista.c
$ ./ricorsione
10 -> 20 -> 30 -> 40 -> NULL
lunghezza: 4
al contrario: 40 30 20 10
\end{console}

\subsection{Come lavorano}

\begin{itemize}
  \item \code{stampaRovescia} stampa la lista al contrario senza mai tornare
        indietro: ogni chiamata \textbf{prima} fa stampare tutto il resto
        della lista, \textbf{poi} stampa il proprio nodo. Il primo a essere
        stampato è quindi il nodo della chiamata più profonda, l'ultimo. Con
        un ciclo sarebbe molto più complicato, perché da un nodo non si può
        tornare al precedente: qui sono le chiamate aperte sulla pila a
        ricordare la strada.
  \item \code{inserisciCodaRic} inserisce in coda un nodo già creato: se la
        lista è vuota, il nodo è tutta la lista; altrimenti lo inserisce in
        coda al resto, e aggancia il risultato a \code{head->next}. È la
        funzione \code{inserisciLista} del progetto delle lezioni.
  \item \code{liberaListaRic} libera prima il resto della lista e poi il
        proprio nodo: così il campo \code{next} viene letto quando il nodo
        esiste ancora.
\end{itemize}

\subsection{Ricorsione o ciclo?}

Tutto ciò che si fa con la ricorsione si può fare con un ciclo, e viceversa.
La ricorsione rende spesso il codice più corto e più vicino alla definizione
del problema; ma ogni chiamata aperta occupa spazio sulla pila, e con una
lista molto lunga le chiamate aperte tutte insieme possono esaurirla. Il
programma \code{profonda.c} crea una lista di \code{n} nodi e ne calcola la
lunghezza in tutti e due i modi:

\begin{affianco}[0.55][c]
\begin{lstlisting}[style=c, name=profonda.c]
#include <stdio.h>
#include "lista.h"

int lunghezzaRic(struct Node *head)
{
   if (head == NULL)
      return 0;
   return 1 + lunghezzaRic(head->next);
}

int main(void)
{
   struct Node *head = NULL;
   int i, n;

   if (scanf("%d", &n) != 1)
      return 1;
   for (i = 0; i < n; i++)
      head = inserisciTesta(head, i);
   printf("con un ciclo: %d\n", lunghezza(head));
   printf("con la ricorsione: %d\n", lunghezzaRic(head));
   liberaLista(head);
   return 0;
}
\end{lstlisting}
\risultato
\begin{console}
$ echo 100000 | ./profonda
con un ciclo: 100000
con la ricorsione: 100000
$ echo 1000000 | ./profonda
con un ciclo: 1000000
Segmentation fault
\end{console}
\wtnota{\code{echo} scrive il numero, e la barra \code{|} lo manda
all'input del programma, come se lo si digitasse.}
\end{affianco}

Con un milione di nodi servono un milione di chiamate aperte insieme, e gli
8 MB che Linux assegna di solito alla pila (capitolo~\ref{chap:stringhe})
non bastano: il programma viene fermato. Il ciclo, invece, usa sempre le
stesse due variabili. Per questo, sulle liste che possono diventare molto
lunghe, si preferiscono i cicli.

\section{I file}

\subsection{Aprire e chiudere un file}

Finora i programmi hanno letto dalla tastiera e scritto sullo schermo, e
hanno usato dei file solo grazie alla redirezione \code{<} e \code{>} del
terminale (capitolo~\ref{chap:primi-programmi}). Il progetto delle lezioni,
invece, apre i file direttamente. Gli strumenti sono in \code{stdio.h}:

\begin{itemize}
  \item \code{FILE *fp;} dichiara un puntatore a \code{FILE}, una struttura
        definita in \code{stdio.h} che descrive un file aperto: quale file
        è, a che punto della lettura o della scrittura si è arrivati, il
        buffer, gli eventuali errori. Non serve sapere com'è fatta: la si usa
        solo attraverso le funzioni della libreria.
  \item \code{fp = fopen("dati.txt", "r");} apre il file \code{dati.txt} in
        \textbf{lettura} (\code{"r"}, \textit{read}). Con \code{"w"}
        (\textit{write}) lo apre in \textbf{scrittura}, creandolo se non
        esiste e \textbf{svuotandolo} se esiste; con \code{"a"}
        (\textit{append}) scrive in fondo a ciò che contiene già. Se il file
        non si può aprire (per esempio non esiste, o manca il permesso),
        \code{fopen} restituisce \code{NULL}, e questo va sempre controllato.
  \item \code{fclose(fp);} chiude il file quando non serve più: scrive su
        disco ciò che è rimasto nel buffer e libera la struttura \code{FILE}.
\end{itemize}

Un nome come \code{"dati.txt"}, senza percorso, indica un file nella cartella
in cui ci si trova quando si lancia il programma: in VS Code, quella del
terminale (capitolo~\ref{chap:ambiente}).

\subsection{Leggere e scrivere}

\code{fscanf} e \code{fprintf} funzionano esattamente come \code{scanf} e
\code{printf}, con un primo argomento in più: il file da cui leggere o su cui
scrivere. Anche il valore restituito da \code{fscanf} è quello di
\code{scanf}: il numero di conversioni riuscite, oppure \code{EOF} se il file
è finito prima della prima conversione. Del resto \code{printf(...)} equivale
a \code{fprintf(stdout, ...)}, e \code{scanf(...)} a
\code{fscanf(stdin, ...)}: \code{stdin} e \code{stdout} sono due
\code{FILE *} già aperti all'avvio del programma, collegati alla tastiera e
allo schermo (o ai file indicati con \code{<} e \code{>}). Il programma
seguente legge dei numeri da \code{numeri.txt} e ne scrive i quadrati in
\code{quadrati.txt}:

\begin{affianco}[0.55][c]
\begin{lstlisting}[style=c, name=quadrati.c]
#include <stdio.h>

int main(void)
{
   FILE *in, *out;
   int x, letti = 0;

   in = fopen("numeri.txt", "r");
   if (in == NULL) {
      printf("non riesco ad aprire numeri.txt\n");
      return 1;
   }
   out = fopen("quadrati.txt", "w");
   if (out == NULL) {
      printf("non riesco a creare quadrati.txt\n");
      fclose(in);
      return 1;
   }
   while (fscanf(in, "%d", &x) == 1) {
      fprintf(out, "%d al quadrato fa %d\n", x, x * x);
      letti++;
   }
   fclose(in);
   fclose(out);
   printf("letti %d numeri\n", letti);
   return 0;
}
\end{lstlisting}
\risultato
\begin{console}
$ cat numeri.txt
3 7
12
$ gcc -Wall -o quadrati quadrati.c
$ ./quadrati
letti 3 numeri
$ cat quadrati.txt
3 al quadrato fa 9
7 al quadrato fa 49
12 al quadrato fa 144
$ rm numeri.txt
$ ./quadrati
non riesco ad aprire numeri.txt
\end{console}
\end{affianco}

\begin{itemize}
  \item Il ciclo \code{while (fscanf(in, "\%d", \&x) == 1)} legge un numero
        alla volta, finché la lettura riesce. \code{fscanf} salta spazi e a
        capo, quindi non importa come i numeri siano disposti nelle righe.
        Alla fine del file restituisce \code{EOF}, e il ciclo termina.
  \item Il comando \code{cat} stampa il contenuto di un file, \code{rm} lo
        cancella. Senza \code{numeri.txt}, \code{fopen} restituisce
        \code{NULL}, e il programma lo dice invece di usare un puntatore
        nullo.
  \item Se la creazione di \code{quadrati.txt} fallisce, il programma chiude
        \code{numeri.txt} prima di uscire: ogni file aperto va chiuso.
\end{itemize}

\section{Il progetto delle lezioni}

\subsection{Che cosa fa}

Torniamo al progetto del capitolo~\ref{chap:make}, che ora abbiamo tutti gli
strumenti per capire. Il file \code{dati.txt} contiene gli esami sostenuti da
alcuni studenti, uno per riga: nome, cognome, matricola, esame e voto (in
fondo al file ci sono anche tre righe vuote, che la lettura salta).

\begin{affianco}[0.55][c]
\begin{lstlisting}[style=plain, name=dati.txt]
Luca      Verdi     364     AnalisiI            16
Enzo      Blu       367     analisiI            14
Mario     Rossi     272     programmazioneI     27
Luca      Verdi     364     basididati1         16
Enzo      Blu       367     basididati1         14
Anna      Bianco    471     algebra             23
Mario     Rossi     272     algebra             30
Luca      Verdi     364     geometria           16
Enzo      Blu       367     algebra             14
\end{lstlisting}
\risultato
Il programma:
\begin{passi}
  \item legge il file e costruisce una lista, con un nodo per ogni esame
        nell'ordine del file, e la stampa;
  \item toglie dalla lista gli esami con voto minore di 18, e la stampa;
  \item scrive nel file \code{dati2.txt} gli esami rimasti, raggruppati per
        studente, e stampa di nuovo la lista;
  \item libera la memoria.
\end{passi}
\end{affianco}

\subsection{\code{studente.h}}

\begin{affianco}[0.55][c]
\codicepiccolo
\begin{lstlisting}[style=h, name=studente.h]
#include <stdio.h>
#include <stdlib.h>
#include <string.h>

#define STRING_LENGTH  512

struct studente{
  char nome[STRING_LENGTH];
  char cognome[STRING_LENGTH];
  int matricola;
  char esame[STRING_LENGTH];
  int voto;
  struct studente *next;
};

void gestioneErrore();
struct studente *leggiFile(FILE *fp, struct studente *lista);
void scriviLista(struct studente *lista);
struct studente *rimuoviSoglia(struct studente *lista, int soglia);
void scriviFile(FILE *fp, struct studente *lista);
void freeLista(struct studente *lista);
\end{lstlisting}
\risultato
\begin{itemize}
  \item Il nodo della lista è una \code{struct studente}: i dati di un esame
        (cinque campi) e il puntatore \code{next}, come in
        \code{struct Node}.
  \item Nome, cognome ed esame sono array di \code{STRING\_LENGTH}
        caratteri dentro il nodo, che occupa così 1552 byte
        (\code{sizeof}), anche se i nomi sono corti.
  \item L'header include \code{stdio.h} perché i prototipi usano il tipo
        \code{FILE}. Sono i prototipi delle funzioni chiamate da
        \code{main.c}; le altre funzioni di \code{funzioni.c} sono
        \code{static} (capitolo~\ref{chap:funzioni}), visibili solo dentro
        quel file.
  \item \code{gestioneErrore()} ha le parentesi vuote, che fino al C17
        vogliono dire \textit{parametri non specificati}: la forma precisa
        è \code{gestioneErrore(void)}.
\end{itemize}
\end{affianco}

\subsection{Creare, inserire e togliere un nodo}

Le prime funzioni di \code{funzioni.c} lavorano su un nodo alla volta:

\begin{lstlisting}[style=c, name=funzioni.c]
#include <stdio.h>
#include <stdlib.h>
#include <string.h>
#include "studente.h"

static struct studente *nuovoNodo(char *nome,char *cognome, int matricola, char *esame, int voto)
{
  struct studente *nodo;

  nodo = (struct studente *) malloc(sizeof(struct studente));
  if (nodo == NULL) return NULL;

  strcpy(nodo->nome, nome);
  strcpy(nodo->cognome, cognome);
  strcpy(nodo->esame, esame);
  nodo->matricola = matricola;
  nodo->voto = voto;
  nodo->next = NULL;

  return nodo;
}

static void freeNodo(struct studente *nodo)
{
  free((void *) nodo);
}

static struct studente *inserisciLista(struct studente *nodo, struct studente *lista)
{
  if (lista==NULL) return nodo;
  lista->next = inserisciLista(nodo, lista->next);
  return lista;
}

static struct studente *rimuoviNodo(struct studente *lista, struct studente *nodo)
{
  struct studente *tmp;


  if (lista == nodo) {
    tmp = lista->next;
    freeNodo(lista);
    return tmp;
  }

  lista->next = rimuoviNodo(lista->next,nodo);
}
\end{lstlisting}

\begin{itemize}
  \item \code{nuovoNodo} è come quella del modulo \code{lista.c}, ma i dati
        sono cinque, e le stringhe vanno \textbf{copiate} nel nodo con
        \code{strcpy} (capitolo~\ref{chap:stringhe}): gli array del
        chiamante verranno riusati per la riga successiva del file. Il cast
        \code{(struct studente *)} davanti a \code{malloc} non è necessario
        in C, come quello a \code{void *} in \code{freeNodo}.
  \item \code{inserisciLista} inserisce \code{nodo} in coda alla lista,
        ricorsivamente: è \code{inserisciCodaRic}, con gli argomenti in
        ordine inverso.
  \item \code{rimuoviNodo} toglie \code{nodo} dalla lista, ricorsivamente:
        se il primo nodo è quello da togliere, lo libera e restituisce il
        resto; altrimenti lo toglie dal resto, e aggancia il risultato a
        \code{lista->next}.
\end{itemize}

Compilando con \code{-Wall}, gcc segnala un problema proprio in
\code{rimuoviNodo}:

\begin{console}
$ gcc -Wall -c main.c funzioni.c
funzioni.c: In function (*@‘@*)rimuoviNodo(*@’@*):
funzioni.c:47:1: warning: control reaches end of non-void function [-Wreturn-type]
   47 | }
      | ^
\end{console}

\warning{Una svista nel codice delle lezioni}{
  Nel ramo ricorsivo di \code{rimuoviNodo} manca il \code{return}: la
  funzione aggancia il risultato a \code{lista->next}, poi arriva alla
  graffa finale senza restituire niente, e il chiamante riceve un valore
  indefinito (in pratica, ciò che si trova per caso nel registro del
  processore usato per i risultati: capitolo~\ref{chap:funzioni}). Nel
  programma l'errore non si manifesta, perché \code{rimuoviNodo} viene
  chiamata solo da \code{rimuoviSoglia}, sempre con \code{nodo} uguale al
  primo nodo della lista: il ramo ricorsivo non viene mai eseguito. Ma
  basterebbe usare la funzione in un altro modo per perdere la lista. La
  correzione è una riga, \code{return lista;} prima della graffa finale:
  con questa, gcc non segnala più niente, e l'output del programma è lo
  stesso.}

\subsection{Le funzioni di servizio}

\begin{lstlisting}[style=c, name=funzioni.c, firstnumber=50]
static int lengthLista(struct studente *lista)
{
  int n = 0;

  while(lista!=NULL) {
    n++;
    lista = lista->next;
  }

  return n;
}


void gestioneErrore()
{
  printf("Errore\n");

  exit(0);
}


void freeLista(struct studente *lista)
{
  if (lista != NULL) {
    freeLista(lista->next);
    freeNodo(lista);
  }
}


void scriviLista(struct studente *lista)
{
  while (lista!=NULL) {
    printf("%s %s %d %s %d\n",lista->nome,lista->cognome,lista->matricola,lista->esame,lista->voto);
    lista = lista->next;
  }
  printf("\n");
}
\end{lstlisting}

\begin{itemize}
  \item \code{lengthLista} conta i nodi con un ciclo, come
        \code{lunghezza}.
  \item \code{gestioneErrore} stampa un messaggio e termina il programma con
        \code{exit} (capitolo~\ref{chap:primi-programmi}), da qualunque
        punto venga chiamata. Termina però con \code{exit(0)}, che alla
        shell comunica un successo: per un errore sarebbe più corretto
        \code{exit(1)}.
  \item \code{freeLista} libera la lista ricorsivamente, prima il resto e poi
        il nodo, come \code{liberaListaRic}.
  \item \code{scriviLista} stampa un nodo per riga, e una riga vuota alla
        fine.
\end{itemize}

\subsection{Leggere il file: \code{leggiFile}}

\begingroup
\codicepiccolo
% numeri di riga a tre cifre: una colonna dei numeri un po' piu' larga
\lstset{xleftmargin=19pt, framexleftmargin=15pt}
\begin{lstlisting}[style=c, name=funzioni.c, firstnumber=89]
struct studente *leggiFile(FILE *fp, struct studente *lista)
{
  char nome[STRING_LENGTH], cognome[STRING_LENGTH], esame[STRING_LENGTH];
  int matricola, voto;
  struct studente *nodo;

  while(fscanf(fp,"%s%s%d%s%d",nome,cognome,&matricola,esame,&voto)==5){
    nodo = nuovoNodo(nome,cognome,matricola,esame,voto);
    if (nodo==NULL) gestioneErrore();
    lista = inserisciLista(nodo, lista);
  }
  return lista;
}
\end{lstlisting}
\endgroup

\code{leggiFile} riceve un file già aperto e la lista a cui aggiungere i
nodi, all'inizio vuota. Il ciclo legge una riga alla volta: \code{fscanf}
legge nome, cognome, matricola, esame e voto, e restituisce 5 se tutte le
conversioni sono riuscite. Per ogni riga si crea un nodo e lo si inserisce in
coda, così la lista ha lo stesso ordine del file; se manca la memoria, si
chiama \code{gestioneErrore}. Quando il file finisce, \code{fscanf}
restituisce \code{EOF}, il ciclo si ferma, e la funzione restituisce la
lista. Il \code{\%s} legge una parola, cioè si ferma al primo spazio: per
questo nel file gli esami sono scritti senza spazi, come
\code{programmazioneI}.

\subsection{Filtrare: \code{rimuoviSoglia}}

\begingroup
\codicepiccolo
% numeri di riga a tre cifre: una colonna dei numeri un po' piu' larga
\lstset{xleftmargin=19pt, framexleftmargin=15pt}
\begin{lstlisting}[style=c, name=funzioni.c, firstnumber=104]
struct studente *rimuoviSoglia(struct studente *lista, int soglia)
{
  if (lista==NULL) return lista;
  if (lista->voto<soglia) {
    lista = rimuoviNodo(lista, lista);
    lista = rimuoviSoglia(lista, soglia);
  } else lista->next = rimuoviSoglia(lista->next, soglia);

  return lista;
}
\end{lstlisting}
\endgroup

\code{rimuoviSoglia} toglie dalla lista tutti i nodi con voto minore di
\code{soglia}, ricorsivamente:

\begin{itemize}
  \item se la lista è vuota, non c'è niente da togliere;
  \item se il primo nodo ha un voto insufficiente, lo toglie
        (\code{rimuoviNodo(lista, lista)} lo libera e restituisce il resto),
        e poi filtra il resto, che è diventato la lista;
  \item altrimenti il primo nodo resta, e la funzione filtra il resto,
        agganciando il risultato a \code{lista->next}.
\end{itemize}

\begin{figure}[H]
  \centering
  \begin{tikzpicture}[font=\Lato\scriptsize, stretto/.style={minimum width=0.62cm}]
    \node[mnota, anchor=west, font=\Lato\scriptsize\bfseries] at (-1.0,0.8)
      {la lista letta da \code{dati.txt}};
    \variabile[minimum width=0.6cm]{h}{(0,0)}{lista}{}
    \nodo[stretto, mlvia]{s1}{(0.75,0)}{16}
    \nodo[stretto, mlvia]{s2}{(2.3,0)}{14}
    \nodo[stretto]{s3}{(3.85,0)}{27}
    \nodo[stretto, mlvia]{s4}{(5.4,0)}{16}
    \nodo[stretto, mlvia]{s5}{(6.95,0)}{14}
    \nodo[stretto]{s6}{(8.5,0)}{23}
    \nodo[stretto]{s7}{(10.05,0)}{30}
    \nodo[stretto, mlvia]{s8}{(11.6,0)}{16}
    \nodo[stretto, mlvia]{s9}{(13.15,0)}{14}
    \nodonull{s9}
    \riferimento{h}{-- (s1.west)}
    \foreach \k/\j in {1/2, 2/3, 3/4, 4/5, 5/6, 6/7, 7/8, 8/9} {\collega{s\k}{s\j}}
    \foreach \k/\nome in {1/Luca, 2/Enzo, 3/Mario, 4/Luca, 5/Enzo, 6/Anna, 7/Mario, 8/Luca, 9/Enzo}
      \node[mindir, anchor=south] at ([yshift=1pt]s\k.north) {\nome};
    \node[mnota, anchor=west, font=\Lato\scriptsize\bfseries] at (-1.0,-0.95)
      {dopo \code{lista = rimuoviSoglia(lista, 18)}};
    \variabile[minimum width=0.6cm]{h2}{(0,-1.75)}{lista}{}
    \nodo[stretto]{r1}{(0.75,-1.75)}{27}
    \nodo[stretto]{r2}{(2.3,-1.75)}{23}
    \nodo[stretto]{r3}{(3.85,-1.75)}{30}
    \nodonull{r3}
    \riferimento{h2}{-- (r1.west)}
    \collega{r1}{r2}
    \collega{r2}{r3}
    \foreach \k/\nome in {1/Mario, 2/Anna, 3/Mario}
      \node[mindir, anchor=south] at ([yshift=1pt]r\k.north) {\nome};
  \end{tikzpicture}
  \caption{I voti dei nove esami letti da \code{dati.txt}, con il nome dello
           studente sopra ogni nodo: \code{rimuoviSoglia} libera i sei nodi
           con voto minore di 18, evidenziati, e collega fra loro quelli che
           restano.}
  \label{fig:soglia}
\end{figure}

\subsection{Raggruppare per studente: \code{scriviFile}}

\begingroup
\codicepiccolo
% numeri di riga a tre cifre: una colonna dei numeri un po' piu' larga
\lstset{xleftmargin=19pt, framexleftmargin=15pt}
\begin{lstlisting}[style=c, name=funzioni.c, firstnumber=116]
void scriviFile(FILE *fp, struct studente *lista)
{
  int *buffer, n, ind1, ind2;
  struct studente *ptr1, *ptr2;

  n = lengthLista(lista);
  buffer = (int *) calloc(n,sizeof(int));

  for (ptr1=lista, ind1=0; ptr1!=NULL; ptr1=ptr1->next, ind1++) {
    if (!buffer[ind1]) {
      buffer[ind1] = 1;
      fprintf(fp,"%s %s %d %s %d\n",ptr1->nome, ptr1->cognome, ptr1->matricola, ptr1->esame, ptr1->voto);
    } else continue;
    for (ptr2=ptr1->next, ind2=ind1+1; ptr2!=NULL; ptr2=ptr2->next, ind2++) {
      if (ptr2->matricola==ptr1->matricola) {
        buffer[ind2] = 1;
        fprintf(fp,"%s %s %d %s %d\n",ptr2->nome, ptr2->cognome, ptr2->matricola, ptr2->esame, ptr2->voto);
      }
    }
  }

  free((void *) buffer);
}
\end{lstlisting}
\endgroup

\code{scriviFile} scrive i nodi nel file \textbf{raggruppati per matricola}:
tutti gli esami di uno studente uno dopo l'altro, nell'ordine in cui
compaiono nella lista.

\begin{itemize}
  \item \code{buffer} è un array di \code{n} interi, uno per nodo, allocato
        con \code{calloc} (capitolo~\ref{chap:puntatori}), quindi tutti a
        zero: \code{buffer[i]} diventa 1 quando il nodo in posizione
        \code{i} è stato scritto. Una lista non ha indici, quindi i due
        \code{for} fanno avanzare insieme due variabili, separate dalla
        virgola: il puntatore al nodo e la sua posizione (\code{ptr1} e
        \code{ind1}, \code{ptr2} e \code{ind2}).
  \item Il \code{for} esterno prende i nodi uno alla volta. Se il nodo è già
        stato scritto, passa al successivo (\code{continue}); altrimenti lo
        scrive, e il \code{for} interno scorre i nodi \textbf{seguenti},
        scrivendo (e segnando) quelli con la stessa matricola.
\end{itemize}

Seguiamo a mano la scrittura della lista filtrata, che ha tre nodi:

\begin{center}
\renewcommand{\arraystretch}{1.2}
\begin{tabular}{@{}l c p{8.4cm}@{}}
  \toprule
  \textbf{\code{ptr1} (\code{ind1})} & \textbf{\code{buffer}} & \textbf{che cosa succede} \\
  \midrule
  Mario Rossi, programmazioneI (0) & \code{0 0 0} & lo scrive; il \code{for}
    interno trova in posizione 2 l'altro esame della matricola 272 e scrive
    anche quello: \code{buffer} diventa \code{1 0 1} \\
  Anna Bianco, algebra (1) & \code{1 0 1} & lo scrive; nessun nodo
    successivo ha la matricola 471: \code{buffer} diventa \code{1 1 1} \\
  Mario Rossi, algebra (2) & \code{1 1 1} & è già stato scritto:
    \code{continue} \\
  \bottomrule
\end{tabular}
\end{center}

\subsection{\code{main.c} e l'esecuzione}

\begin{affianco}[0.5][c]
\begin{lstlisting}[style=c, name=main.c]
#include <stdio.h>
#include <stdlib.h>
#include <string.h>
#include "studente.h"

int main()
{
  struct studente *lista=NULL;
  FILE *fp;

  fp = fopen("dati.txt","r");
  if (fp==NULL) gestioneErrore();
  lista = leggiFile(fp, lista);
  fclose(fp);
  scriviLista(lista);

  lista = rimuoviSoglia(lista, 18);
  scriviLista(lista);

  fp = fopen("dati2.txt","w");
  if (fp==NULL) gestioneErrore();
  scriviFile(fp, lista);
  fclose(fp);

  scriviLista(lista);

  freeLista(lista);
}
\end{lstlisting}
\risultato
\begin{console}
$ gcc -Wall -c main.c funzioni.c
(*@\textit{\dots{} l'avviso su rimuoviNodo \dots}@*)
$ gcc -o esercizio main.o funzioni.o
$ ./esercizio
Luca Verdi 364 AnalisiI 16
Enzo Blu 367 analisiI 14
Mario Rossi 272 programmazioneI 27
Luca Verdi 364 basididati1 16
Enzo Blu 367 basididati1 14
Anna Bianco 471 algebra 23
Mario Rossi 272 algebra 30
Luca Verdi 364 geometria 16
Enzo Blu 367 algebra 14

Mario Rossi 272 programmazioneI 27
Anna Bianco 471 algebra 23
Mario Rossi 272 algebra 30

Mario Rossi 272 programmazioneI 27
Anna Bianco 471 algebra 23
Mario Rossi 272 algebra 30

$ cat dati2.txt
Mario Rossi 272 programmazioneI 27
Mario Rossi 272 algebra 30
Anna Bianco 471 algebra 23
\end{console}
\end{affianco}

\begin{itemize}
  \item \code{main} apre \code{dati.txt} in lettura (se non ci riesce chiama
        \code{gestioneErrore}), costruisce la lista con \code{leggiFile},
        chiude il file e stampa i nove esami.
  \item \code{rimuoviSoglia(lista, 18)} lascia i tre esami superati, e la
        lista viene stampata di nuovo.
  \item \code{dati2.txt} viene aperto in scrittura e riempito da
        \code{scriviFile}: i due esami di Mario Rossi sono uno dopo l'altro,
        come voleva il raggruppamento. La terza stampa è uguale alla
        seconda, perché \code{scriviFile} non modifica la lista.
  \item Alla fine \code{freeLista} libera i nodi. \code{main} non ha un
        \code{return}: dal C99, arrivare alla graffa finale di \code{main}
        equivale a \code{return 0}.
\end{itemize}

\subsection{Nessuna perdita di memoria}

\textbf{valgrind} è un programma che ne esegue un altro tenendo traccia di
ogni allocazione e di ogni \code{free}, e alla fine dice se qualche blocco è
rimasto allocato: una \textbf{perdita di memoria} (\textit{memory leak}). Sul
progetto (qui l'output del programma, che si mescola a quello di valgrind, è
omesso):

\begin{console}
$ valgrind ./esercizio
==451== Memcheck, a memory error detector
==451== Copyright (C) 2002-2022, and GNU GPL'd, by Julian Seward et al.
==451== Using Valgrind-3.22.0 and LibVEX; rerun with -h for copyright info
==451== Command: ./esercizio
==451==
(*@\textit{\dots{} l'output del programma \dots}@*)
==451==
==451== HEAP SUMMARY:
==451==     in use at exit: 0 bytes in 0 blocks
==451==   total heap usage: 15 allocs, 15 frees, 27,212 bytes allocated
==451==
==451== All heap blocks were freed -- no leaks are possible
==451==
==451== For lists of detected and suppressed errors, rerun with: -s
==451== ERROR SUMMARY: 0 errors from 0 contexts (suppressed: 0 from 0)
\end{console}

Le 15 allocazioni comprendono i 9 nodi, il \code{buffer} di
\code{scriviFile} e quelle che la libreria fa per gestire i file e l'output:
sono state liberate tutte. Proviamo a commentare la chiamata
\code{freeLista(lista);} in \code{main.c}, e a compilare con \code{-g}
(che aggiunge i numeri di riga, utili a strumenti come questo); con
\code{--leak-check=full}, valgrind dice anche dove sono stati allocati i
blocchi persi:

\begin{console}
$ gcc -g -o esercizio main.c funzioni.c
$ valgrind --leak-check=full ./esercizio
(*@\textit{\dots}@*)
==442== HEAP SUMMARY:
==442==     in use at exit: 4,656 bytes in 3 blocks
==442==   total heap usage: 15 allocs, 12 frees, 27,212 bytes allocated
==442==
==442== 4,656 (1,552 direct, 3,104 indirect) bytes in 1 blocks are definitely lost in loss record 2 of 2
==442==    at 0x4846828: malloc (in /usr/libexec/valgrind/vgpreload_memcheck-amd64-linux.so)
==442==    by 0x1093E3: nuovoNodo (funzioni.c:10)
==442==    by 0x1096A9: leggiFile (funzioni.c:96)
==442==    by 0x10931D: main (main.c:13)
==442==
==442== LEAK SUMMARY:
==442==    definitely lost: 1,552 bytes in 1 blocks
==442==    indirectly lost: 3,104 bytes in 2 blocks
(*@\textit{\dots}@*)
\end{console}

I tre nodi rimasti dopo il filtro non sono stati liberati: 1552 byte
ciascuno, come diceva \code{sizeof}. Il primo è \textit{definitely lost}
(perso di sicuro): nessun puntatore lo raggiunge più, perché la variabile
\code{lista} di \code{main} non esiste più. Gli altri due sono
\textit{indirectly lost} (persi indirettamente): sono raggiungibili solo
attraverso il primo, che è perso. È la struttura della lista vista da
valgrind.

\info{Installare valgrind}{
  Su Ubuntu e su WSL valgrind si installa con
  \code{sudo apt install valgrind}, e si usa scrivendo \code{valgrind} prima
  del comando che lancia il programma. Su macOS valgrind non funziona: per
  gli accessi sbagliati alla memoria c'è \code{-fsanitize=address}, e per le
  perdite il comando \code{leaks} (\code{leaks --atExit -- ./esercizio}).}

\insintesi{
  Una \textbf{lista collegata} è una catena di nodi allocati nello heap: ogni
  nodo contiene i dati e l'indirizzo del successivo, l'ultimo ha \code{next}
  uguale a \code{NULL}, e la lista è individuata dal puntatore al primo nodo
  (\code{NULL} se è vuota). Rispetto a un array, la dimensione è flessibile e
  inserire o togliere un nodo costa il cambio di un paio di puntatori, ma
  l'accesso è sequenziale. Quasi ogni operazione segue lo stesso schema: un
  puntatore di appoggio scorre la lista e si ferma sul nodo giusto, spesso il
  \textbf{precedente} di quello su cui si opera; l'ordine delle assegnazioni
  conta, e un nodo si libera solo dopo averne letto il \code{next}. Le
  funzioni che possono cambiare il primo nodo restituiscono la nuova testa.
  Una lista è anche una struttura \textbf{ricorsiva} (vuota, oppure un nodo
  seguito da una lista), e molte operazioni si scrivono con funzioni che
  chiamano se stesse, con un caso base e un passo ricorsivo; ogni chiamata
  aperta, però, occupa spazio sulla pila. I file si aprono con \code{fopen},
  si usano con \code{fscanf} e \code{fprintf} e si chiudono con
  \code{fclose}. valgrind e \code{-fsanitize=address} aiutano a scoprire le
  perdite di memoria e gli accessi alla memoria già liberata.}

\section{Riferimenti}

\begin{itemize}
  \item \textbf{A. Kelley, I. Pohl}, \textit{A Book on C}: capitolo~10
        (strutture e liste), capitolo~11 (input/output) e, per la
        ricorsione, capitolo~5 (funzioni).
  \item \textbf{fopen} --- cppreference.com. \\
        \urlx{https://en.cppreference.com/w/c/io/fopen}
  \item \textbf{Valgrind Quick Start Guide} --- valgrind.org. \\
        \urlx{https://valgrind.org/docs/manual/quick-start.html}
\end{itemize}
